% % % % % % % % % % % % % % % % % % % % % % % % % % %
% IS&T Template 
% Patrick Vandewalle
% January 2006
% % % % % % % % % % % % % % % % % % % % % % % % % % %

%%%%%%%%%%%%%%%%%%%%%%%%%%%%%%%%%%
% Document class
%%%%%%%%%%%%%%%%%%%%%%%%%%%%%%%%%%
\documentclass[letterpaper,twocolumn,fleqn]{article} 

%%%%%%%%%%%%%%%%%%%%%%%%%%%%%%%%%%
% Packages
%%%%%%%%%%%%%%%%%%%%%%%%%%%%%%%%%%
\usepackage{ist}
\usepackage[margin=1in]{geometry}
% add other packages here
\usepackage{float} %forza
\usepackage{wrapfig} %forza
\usepackage{amsmath}
\usepackage{amssymb}
\usepackage{graphicx}
\usepackage{listings}
\usepackage{xcolor}
\usepackage{float}
\usepackage{hyperref}
\usepackage{booktabs}
\usepackage[american]{circuitikz}
\pagestyle{empty}                % no page numbers is default
\lstset{
  basicstyle=\ttfamily\footnotesize,
  breaklines=true,
  breakatwhitespace=true,
  columns=fullflexible,
  keepspaces=true,
  frame=single,
  showstringspaces=false,
  xleftmargin=0.02\columnwidth,
  xrightmargin=0.02\columnwidth
}

%%%%%%%%%%%%%%%%%%%%%%%%%%%%%%%%%%
% Title and Authors
%%%%%%%%%%%%%%%%%%%%%%%%%%%%%%%%%%

\title{Unit 2: Polynomial Interpolation Applied to\\
  Diode I--V Characteristic Reconstruction }
\author{ Jacky Li; Cal Poly Pomona; Pomona, California \\ February 23, 2025}
 % date has an empty field.
% correct for bad hyphenation here
% \hyphenation{}

%%%%%%%%%%%%%%%%%%%%%%%%%%%%%%%%%%
% Begin document
%%%%%%%%%%%%%%%%%%%%%%%%%%%%%%%%%%
\begin{document} 
\maketitle 
\thispagestyle{empty} % prevents the first page to be numbered

%%%%%%%%%%%%%%%%%%%%%%%%%%%%%%%%%%
% Abstract
%%%%%%%%%%%%%%%%%%%%%%%%%%%%%%%%%%

\begin{abstract}
In many engineering scenarios, we have a limited set of measured data points and
need to reconstruct the underlying continuous function.  Polynomial
interpolation provides a systematic way to achieve this.  In this report, we
apply two interpolation methods---\textbf{Lagrange interpolation} and
\textbf{Newton's divided-difference interpolation}---to reconstruct the I--V
characteristic of a silicon diode from simulated lab measurements.  We show that
both methods produce the \emph{same} interpolating polynomial but differ in
computational structure: Lagrange uses basis polynomials while Newton uses a
divided-difference table.  We compare their results against the known Shockley
model and analyze the interpolation error.
\end{abstract}


\section{Introduction}

Polynomial interpolation is the process of constructing a polynomial of degree
at most $n-1$ that passes exactly through $n$ given data points
$(x_0, y_0), (x_1, y_1), \ldots, (x_{n-1}, y_{n-1})$.  Two classical forms of
the interpolating polynomial are:

\begin{itemize}
  \item \textbf{Lagrange form} --- expresses the polynomial as a weighted sum of
        basis polynomials $L_i(x)$, each designed to be~1 at~$x_i$ and~0 at all
        other data points.
  \item \textbf{Newton's divided-difference form} --- builds the polynomial
        incrementally using a divided-difference table, making it easy to add
        new data points without recomputing everything.
\end{itemize}

Both forms yield the \emph{same unique polynomial} (by the uniqueness theorem),
but their computational approaches and practical trade-offs differ.

In this report, we apply both methods to a practical EE problem: reconstructing
the I--V characteristic curve of a silicon diode from a small number of measured
data points.

% ════════════════════════════════════════════════════════════════════
\section{Background}

\subsection{Lagrange Interpolation}
Given $n$ data points, the Lagrange interpolating polynomial is
\begin{equation}
  P(x) = \sum_{i=0}^{n-1} y_i \, L_i(x), \quad
  L_i(x) = \prod_{\substack{j=0 \\ j \neq i}}^{n-1}
            \frac{x - x_j}{x_i - x_j}
  \label{eq:lagrange}
\end{equation}
Each $L_i(x)$ is a polynomial of degree $n-1$ satisfying $L_i(x_j) = \delta_{ij}$.

\subsection{Newton's Divided-Difference Interpolation}
The Newton form is
\begin{equation}
  P(x) = c_0 + c_1(x - x_0) + c_2(x - x_0)(x - x_1) + \cdots
  \label{eq:newton}
\end{equation}
where the coefficients $c_k$ are the \emph{divided differences}:
\begin{align}
  f[x_i] &= y_i \\
  f[x_i, x_{i+1}] &= \frac{f[x_{i+1}] - f[x_i]}{x_{i+1} - x_i} \\
  f[x_i, \ldots, x_{i+k}] &= \frac{f[x_{i+1}, \ldots, x_{i+k}]
                                   - f[x_i, \ldots, x_{i+k-1}]}
                                   {x_{i+k} - x_i}
\end{align}
and $c_k = f[x_0, x_1, \ldots, x_k]$.


% ════════════════════════════════════════════════════════════════════
\section{Problem Formulation}

\subsection{Motivation}
A diode's I--V relationship follows the Shockley equation:
\begin{equation}
  I_d = I_s \left( e^{\,V_d / (n\,V_T)} - 1 \right)
  \label{eq:shockley}
\end{equation}
In practice, we often only have \emph{measured} data from the lab rather than
the exact model parameters.  Interpolation allows us to reconstruct the full
curve from a few measurements.

\subsection{Data}
We use five simulated ``measured'' data points from a silicon diode with
parameters $I_s = 10^{-12}\;\text{A}$, $n = 1.5$, and $V_T = 25.85\;\text{mV}$,
with $\pm2\%$ measurement noise added:

\begin{table}[H]
\centering
\begin{tabular}{@{}ccc@{}}
\toprule
\textbf{Point} & $V_d$ [V] & $I_d$ [mA] \\
\midrule
1 & 0.30 & $2.31 \times 10^{-6}$ \\
2 & 0.50 & $3.97 \times 10^{-4}$ \\
3 & 0.65 & $0.0193$ \\
4 & 0.75 & $0.2590$ \\
5 & 0.85 & $3.2982$ \\
\bottomrule
\end{tabular}
\caption{Measured diode data points (with $\pm2\%$ noise).}
\label{tab:data}
\end{table}

% ════════════════════════════════════════════════════════════════════
% ════════════════════════════════════════════════════════════════════
\section{Implementation}

\subsection{Lagrange Interpolation}
\begin{lstlisting}[caption={Lagrange interpolation (tools.py)},label={lst:lagrange}]
def lagrange_interpolate(self, x, y):
    n = len(x)
    P = np.zeros(n)
    for i in range(n):
        Li = np.array([1.0])
        denom = 1.0
        for j in range(n):
            if i != j:
                Li = np.convolve(Li, np.array([1.0, -x[j]]))
                denom *= (x[i] - x[j])
        P = P + Li * (y[i] / denom)
    return P, 0
\end{lstlisting}

\subsection{Newton's Divided-Difference Interpolation}
\begin{lstlisting}[caption={Newton's divided-difference interpolation (tools.py)},label={lst:newton}]
def newton_interpolate(self, x, y):
    n = len(x)
    dd = np.copy(y).astype(float)
    coefs = [dd[0]]
    for j in range(1, n):
        for i in range(n - 1, j - 1, -1):
            dd[i] = (dd[i] - dd[i-1]) / (x[i] - x[i-j])
        coefs.append(dd[j])
    # Convert to standard polynomial
    P = np.array([coefs[0]])
    basis = np.array([1.0])
    for k in range(1, n):
        basis = np.convolve(basis, np.array([1.0, -x[k-1]]))
        P = np.polyadd(P, coefs[k] * basis)
    return P, 0
\end{lstlisting}

% ════════════════════════════════════════════════════════════════════
\section{Results}

\subsection{Polynomial Coefficients}
Both methods produce the \textbf{same} 4th-degree polynomial (up to floating-point
precision), confirming the uniqueness theorem of polynomial interpolation.

\subsection{I--V Curve Reconstruction}
Figure~\ref{fig:iv} (left) shows the interpolated curves overlaid with the true
Shockley model.  Both Lagrange (blue dashed) and Newton (red dotted) traces are
indistinguishable, and they closely follow the true model within the data range
$[0.3, 0.85]$~V.

Figure~\ref{fig:iv} (right) shows the absolute interpolation error on a
logarithmic scale.  The error is smallest near the data points and grows toward
the edges of the interpolation interval.

\begin{figure}[H]
  \centering
  \includegraphics[width=0.5\textwidth]{unit02_interpolation_results.png}
  \caption{Left: Interpolated diode I--V curves vs.\ true Shockley model.
           Right: Absolute error of interpolation vs.\ true model.}
  \label{fig:iv}
\end{figure}

\subsection{Data Trend Comparison}
Figure~\ref{fig:trend} compares the true diode current, Lagrange prediction,
and Newton prediction at several test voltages \emph{not} present in the
original data set.  The close agreement confirms accurate interpolation within
the data range.

\begin{figure}[H]
  \centering
  \includegraphics[width=0.5\textwidth]{unit02_data_trend.png}
  \caption{Data trend: True vs.\ Lagrange vs.\ Newton at test voltages.}
  \label{fig:trend}
\end{figure}

% ════════════════════════════════════════════════════════════════════
\section{Discussion}

\subsection{Lagrange vs.\ Newton: Comparison}
\begin{table}[H]
\centering
\small
\begin{tabular}{@{}lcc@{}}
\toprule
\textbf{Property} & \textbf{Lagrange} & \textbf{Newton} \\
\midrule
Produces unique polynomial & Yes & Yes \\
Easy to add new data points & No (recompute all) & Yes (extend table) \\
Computational form & Basis polynomials & Divided differences \\
Numerical stability & May suffer for large $n$ & Generally better \\
Good for understanding & Yes (explicit formula) & Yes (recursive) \\
\bottomrule
\end{tabular}
\caption{Comparison of interpolation methods.}
\label{tab:comparison}
\end{table}

\subsection{Practical Considerations}
\begin{itemize}
  \item Both methods produce the same result for the same data, but Newton's
        form is more efficient when data points are added incrementally (e.g.,
        during a lab experiment).
  \item With only $n=5$ data points, the 4th-degree polynomial provides a
        reasonable approximation within the data range but should \emph{not}
        be used for extrapolation.
  \item For the exponential nature of the diode curve, higher-order
        interpolation or specialized fitting (e.g., nonlinear least squares)
        may be preferred for high-accuracy applications.
\end{itemize}

% ════════════════════════════════════════════════════════════════════
\section{Conclusion}

We applied Lagrange and Newton's divided-difference interpolation to
reconstruct a silicon diode's I--V characteristic from five measured data
points.  Both methods produced the identical 4th-degree polynomial, verifying
the uniqueness theorem.  The interpolated curve closely matched the true
Shockley model within the data range.  Newton's form offers practical
advantages when incrementally adding data, while Lagrange's form provides a
more intuitive closed-form expression.

\end{document}



\clearpage
%%%%%%%%%%%%%%%%%%%%%%%%%%%%%%%%%%%%
% Overall Document Guidelines: Head
%%%%%%%%%%%%%%%%%%%%%%%%%%%%%%%%%%%%
\section{Overall Document Guidelines: Head}
\label{sec:intro}

Clear document of any fonts other than Times, Arial, and Symbol. The
paper should be formatted using the tags provided in the template,
i.e., title, author, section, subsection, subsubsection, eq./fig.,
references, etc.  in a 2 column format on US letter size paper ($8.5
\times 11$ inches, or $21.6 \times 27.9$ cm).

The left and right margins are set automatically to .75 inch (1.90 cm),
and the top and bottom margins to 1.0 inch (2.54 cm). The document is in
a 2-column format with column widths set at 3.38 inch (8.57 cm) and the
gutter -the space between columns- at .25 inch (0.635 cm).

Papers should be a maximum of 4-6 pages; longer papers will be
returned for revision. Please do not place folios or page numbers in
your paper. That information is inserted when we assemble the book.


%%%%%%%%%%%%%%%%%%%%%%%%%%%%%%%%%%
% Graphics and Equations
%%%%%%%%%%%%%%%%%%%%%%%%%%%%%%%%%%
\section{Graphics and Equations}
Graphics and equations should fit within one column (3.38 inches
wide), but full width (7 inch) figures are also acceptable. Equations,
figures and figure captions each have their own style tags. Equations
are numbered using parentheses flushed right as shown below.

\begin{equation}
\label{eq:ist}
\textrm{IS\&T} + \textrm{members} \times \textrm{Confs.} = \textrm{Success}
\end{equation}


\subsection{Helpful Hints and Style Tags: Subhead}
For a complete listing of the style tags for use in this template
refer to Table \ref{tab:fonts}. These are the style tags for
conference proceedings; if you use the wrong template/style tags your
paper will be sent back to you to be reformatted.  All of these forms
and templates related to the publication of conference papers are
available at

www.imaging.org/conferences/guidelines.cfm. 

Select the specific conference and download the Authors Kit. The
template may vary from one conference to another.

The template is set up for MS Word and LaTeX. Please check the paper
carefully to confirm that the styles have been applied correctly, then
print it out and double check to ensure that the paper appears as
intended.

\begin{table}[!h]
\caption{Style Tag Table: Table head}
\label{tab:fonts}
\begin{center}       
\begin{tabular}{|p{0.45\columnwidth}|p{0.45\columnwidth}|} 
\hline
1. Title& title environment \\ \hline
2. Byline: & \\
Author/Affiliation & author environment \\ \hline
3. Head & section environment \\ \hline
4. Subhead & subsection environment \\ \hline
5. Tertiary head & subsubsection environment \\ \hline
6. Abstract & abstract environment \\ \hline 
& \\ \hline
7. Body & regular text\\ \hline
8. Equation & equation environment\\ \hline
9.1 Figure & align left \\ \hline
9.2 Figure caption & caption environment \\ \hline
10.1 Table head & caption environment \\ \hline
10.2 Table text & tabular environment \\ \hline
11. References & bibliography + 
command \emph{small} to have a 1pt smaller font\\ \hline
12. Author Bio text & biography environment\\ \hline
13.1 List, bullet & itemize environment\\ \hline
13.2 List, numbered & enumerate environment \\ \hline
\end{tabular}
\end{center}
\end{table} 

%%%%%%%%%%%%%%%%%%%%%%%%%%%%%%%%%%
% Submitting Your Paper
%%%%%%%%%%%%%%%%%%%%%%%%%%%%%%%%%%

\section{Submitting Your Paper}
The submission of your paper has to be performed through the IS\&T
submission website. Authors receive a login for this site by e-mail.
Papers can be submitted in Postscript, Word, or WordPerfect format,
and will be converted to PDF by the submission server. Please
carefully review the generated PDF and verify that all the text,
equations, figures and tables are displayed correctly before approving
its submission.

\subsection{Margins in LaTeX}
Because of the differences between dvips conversion utilities, the
margins of your generated PDF document might vary. Please print
your document, and verify its margins. If they are incorrect, please
adapt the sizes of the margins in the file \emph{ist.sty}. Typically the
top margin should be decreased.

\begin{table}[!h]
\caption{Document Specs: Table head}
\label{tab:specs}
\begin{center}       
\begin{tabular}{p{0.35\columnwidth}p{0.55\columnwidth}} 
Paper Size & US Letter \\
Left/right margin & .75 inch (1.90 cm) \\
Top/bottom margin & 1 inch (2.54 cm) \\
Columns & 2 at 3.38 inch (8.57 cm) wide. \\
 & Spacing between columns: 0.25 inch (0.635 cm)
\end{tabular}
\end{center}
\end{table} 

\begin{figure}[!hb]
  \includegraphics[width=0.3\columnwidth]{logo.png}
  \caption{IS\&T logo.}
  \label{Figure:logo}
\end{figure}

Please contact IS\&T with any questions or requests for assistance in
helping prepare the paper. We look forward to having your paper
presented at the conference and published in the conference
proceedings.

\begin{figure}[!hb]
  \includegraphics[width=0.3\columnwidth]{logo.png}
  \caption{IS\&T logo.}
  \label{Figure:logo}
\end{figure}

%%%%%%%%%%%%%%%%%%%%%%%%%%%%%%%%%%
% Reference Preparation
%%%%%%%%%%%%%%%%%%%%%%%%%%%%%%%%%%

\section{Reference Preparation}
Use the standard LaTeX \emph{cite} command for references in the
text. You can then use the standard bibliography command to generate
the list of references. Add the command \emph{small} before the
bibliography to give it the right font size.  Reference \cite{bib1}
style should be used for books, Reference \cite{bib2} style should be
used for Journals, and Reference \cite{bib3} style should be used for
Proceedings.

%\section{Acknowledgments} 
%add the acknowledgement section here

% To start a new column (but not a new page) and help balance the last-page
% column length use \vfill\pagebreak.

%%%%%%%%%%%%%%%%%%%%%%%%%%%%%%%%%%
% Bibliography
%%%%%%%%%%%%%%%%%%%%%%%%%%%%%%%%%%

\small
\begin{thebibliography}{9}
\bibitem{bib1}John Doe, Recent Progress in Digital Halftoning II,
  IS\&T, Springfield, VA, 1999, pg. 173.
\bibitem{bib2}John Doe, Digital Imaging, J. Imaging. Sci. and
  Technol., 42, 112 (1998).
\bibitem{bib3}John Doe, An Inexpensive Micro-Goniophotometry You Can
  Build, Proc. PICS, pg. 179. (1998).
\end{thebibliography}

%%%%%%%%%%%%%%%%%%%%%%%%%%%%%%%%%%
% Biography
%%%%%%%%%%%%%%%%%%%%%%%%%%%%%%%%%%

\begin{biography}
Please submit a brief biographical sketch of no more than 75 words. 
Include relevant professional and educational information as shown 
in the example below.

Jane Doe received her BS in physics from the University of Nevada (1977) 
and her PhD in applied physics from Columbia University (1983). Since 
then she has worked in the Research and Technology Division at Xerox 
in Webster, NY. Her work has focused on the development of toner adhesion 
and transport issues. She is on the Board of  IS\&T and a member of APS 
and SPIE.
\end{biography}

\end{document} 
